\documentclass[main.tex]{subfiles}
\begin{document}
\ifSubfilesClassLoaded{\begin{center}\small\itshape Part of the GnssProcessor Mathematical Reference \textbullet\ \today\end{center}}{}
\section{Coordinate Frames and Time Systems}\label{sec:frames}
% ============================================================
%  The geodetic-frame and time-scale substrate the engine assumes
% ============================================================

% --- Doc-local symbols (descriptive names to avoid package collisions) ---
\newcommand{\semia}{a}                               % WGS84 semi-major axis
\newcommand{\wgsflat}{f}                              % WGS84 flattening
\newcommand{\eccsq}{e^{2}}                            % first eccentricity squared
\newcommand{\primevn}{N}                              % prime-vertical radius of curvature
\newcommand{\geolat}{\varphi}                         % geodetic latitude
\newcommand{\geolon}{\lambda}                         % geodetic longitude
\newcommand{\geoht}{h}                                % ellipsoidal height
\newcommand{\enuE}{\hat{\bm{e}}_{\mathrm{E}}}         % ENU East basis vector
\newcommand{\enuN}{\hat{\bm{e}}_{\mathrm{N}}}         % ENU North basis vector
\newcommand{\enuU}{\hat{\bm{e}}_{\mathrm{U}}}         % ENU Up basis vector
\newcommand{\enuR}{\bm{R}_{\mathrm{ENU}}}             % ENU->ECEF rotation
\newcommand{\eccentricityHEN}{\bm{e}_{\mathrm{HEN}}}  % RINEX ANTENNA: DELTA H/E/N eccentricity
\newcommand{\elev}{\mathrm{el}}                       % elevation angle
\newcommand{\azim}{\mathrm{az}}                       % azimuth angle
\newcommand{\leapgu}{\Delta_{\mathrm{GU}}}            % GPS - UTC leap seconds
\newcommand{\precP}{\bm{P}_{\mathrm{prec}}}           % precession matrix
\newcommand{\nutN}{\bm{N}_{\mathrm{nut}}}             % nutation matrix
\newcommand{\polW}{\bm{W}_{\mathrm{pm}}}              % polar-motion matrix
\newcommand{\Umat}{\bm{U}}                            % full ECI->ECEF matrix
\newcommand{\rotR}[1]{\bm{R}_{#1}}                    % elementary axis rotation
\newcommand{\dpsi}{\Delta\psi}                        % nutation in longitude
\newcommand{\deps}{\Delta\varepsilon}                % nutation in obliquity
\newcommand{\obleps}{\varepsilon}                     % mean obliquity of the ecliptic
\newcommand{\gmst}{\mathrm{GMST}}
\newcommand{\gast}{\mathrm{GAST}}
\newcommand{\era}{\theta_{\mathrm{ERA}}}             % Earth Rotation Angle
\newcommand{\cipX}{X}                                 % CIP coordinate X
\newcommand{\cipY}{Y}                                 % CIP coordinate Y
\newcommand{\cioloc}{s}                               % CIO locator

Every model in this reference --- the tropospheric and ionospheric
mappings of \dcref{sec:atmos}{the Atmosphere document}, the
solid-Earth, ocean, and pole tides of \dcref{sec:sitedisp}{the Site
Displacement document}, the antenna, attitude, and phase-wind-up terms
of \dcref{sec:obsphys}{the Observation Physics document}, and the orbit
evaluation of \dcref{sec:products}{the Product Evaluation document} ---
presumes a fixed geodetic-frame and time-scale substrate.  This document
specifies that substrate: the WGS84 ellipsoid and the three coordinate
representations (Cartesian ECEF, geodetic, and local ENU) the engine
moves between; the family of GNSS time scales together with the leap-
second machinery that ties them to UTC; and the Earth-orientation chain
that rotates the celestial frame --- in which the Sun, Moon, and the
satellite ephemerides are most naturally expressed --- into the
terrestrial frame in which the receiver sits.  None of the conversions
here is an estimator state.  The one frame quantity that \emph{is}
estimated --- the inter-system time offset (ISB) that absorbs the
unmodelled bias between constellation time scales --- belongs to
\dcref{sec:state}{the State Vector document}; this document covers only
the \emph{deterministic} scale conversions that the ISB state is defined
relative to.  Code references are to \texttt{src/core/math/Coordinates.cpp},
\texttt{src/core/math/TimeUtils.cpp}, \texttt{src/core/model/Planetary.cpp},
\texttt{src/core/products/LiveEphemerisStore.cpp},
\texttt{src/core/products/ProductParser.cpp}, and
\texttt{src/core/rinex/RinexParser.cpp}.

% ============================================================
\subsection{The WGS84 ellipsoid: ECEF, geodetic, and local frames}
\label{sec:frame-geodetic}
% ============================================================

The terrestrial reference frame is the WGS84 Earth-Centred,
Earth-Fixed (ECEF) Cartesian system: origin at the geocentre, $z$ along
the conventional pole, $x$ toward the IERS reference meridian, $y$
completing a right-handed triad.  The defining ellipsoid has semi-major
axis $\semia$ and flattening $\wgsflat$, from which the first eccentricity
squared follows (\texttt{Coordinates::ecefToLlh} and
\texttt{Coordinates::llhToEcef} each derive it from $\wgsflat$):
\begin{equation}
\eccsq \;=\; \wgsflat\,(2-\wgsflat).
\label{eq:frame-e2}
\end{equation}
The radius of curvature in the prime vertical at geodetic latitude
$\geolat$ is
\begin{equation}
\primevn(\geolat) \;=\;
\frac{\semia}{\sqrt{1 - \eccsq\,\sin^2\!\geolat}} ,
\label{eq:frame-primevert}
\end{equation}
which scales the entire forward map from geodetic coordinates
$(\geolat,\geolon,\geoht)$ to ECEF
(\texttt{Coordinates::llhToEcef}):
\begin{empheq}[box=\fbox]{equation}
\begin{aligned}
x &= \big(\primevn + \geoht\big)\cos\geolat\,\cos\geolon, \\
y &= \big(\primevn + \geoht\big)\cos\geolat\,\sin\geolon, \\
z &= \big((1-\eccsq)\,\primevn + \geoht\big)\sin\geolat .
\end{aligned}
\label{eq:frame-llh2ecef}
\end{empheq}
The factor $(1-\eccsq)\primevn$ on the $z$ component is the source of all
ellipsoidal subtlety: a sphere would carry $\primevn$ unmodified on every
axis, and the $(1-\eccsq)$ correction is precisely the contraction that
distinguishes the geodetic normal (perpendicular to the ellipsoid
surface) from the geocentric radius vector.

\paragraph{Inverse map (Bowring fixed point).}
The inverse problem --- recovering $(\geolat,\geolon,\geoht)$ from
$(x,y,z)$ --- has no closed form for the latitude because $\primevn$ in
\eqref{eq:frame-llh2ecef} itself depends on $\geolat$.  The longitude is
immediate, $\geolon = \operatorname{atan2}(y,x)$.  For the latitude the
engine runs a Bowring fixed-point iteration on the equatorial projection
$p = \sqrt{x^2+y^2}$ (\texttt{Coordinates::ecefToLlh}).
The seed neglects height,
\begin{equation}
\geolat^{(0)} = \operatorname{atan2}\!\big(z,\; p\,(1-\eccsq)\big),
\label{eq:frame-bowring-seed}
\end{equation}
and each pass recomputes $\primevn$ and the height and re-evaluates the
latitude with the height feedback folded back into the eccentricity term:
\begin{equation}
\geoht^{(k)} = \frac{p}{\cos\geolat^{(k)}} - \primevn^{(k)},
\qquad
\geolat^{(k+1)} = \operatorname{atan2}\!\left(z,\;
p\left[1 - \eccsq\,\frac{\primevn^{(k)}}{\primevn^{(k)}+\geoht^{(k)}}\right]\right).
\label{eq:frame-bowring-iter}
\end{equation}
The iteration is capped at ten passes and exits early once
$|\geolat^{(k+1)}-\geolat^{(k)}| < 10^{-13}$; for terrestrial heights it
converges to machine precision in two or three steps.

\paragraph{Pole branch.}
When $p < 10^{-12}\,\mathrm{m}$ the point lies on (or within a
picometre of) the $z$ axis and both the longitude and the
$p/\cos\geolat$ height formula degenerate.  The code takes an explicit
branch (\texttt{Coordinates::ecefToLlh}, the $p<10^{-12}$ guard): the
longitude retains the $\operatorname{atan2}(y,x)$ value computed before
the guard --- formally undefined on the axis, and immaterial there ---
while the branch fixes $\geolat = \pm\pi/2$ by the sign of $z$ and reads
the height directly off the polar semi-minor axis,
\begin{equation}
\geoht = |z| - \primevn\,(1-\eccsq)
       = |z| - \semia\sqrt{1-\eccsq} \quad (\text{at the pole, } \sin^2\!\geolat = 1),
\label{eq:frame-pole}
\end{equation}
since $\primevn(1-\eccsq)$ evaluated at the pole is exactly the polar
radius $\semia\sqrt{1-\eccsq}$.  This guard keeps the inverse map total
over all of $\mathbb{R}^3$, which matters because the planetary and
attitude code occasionally evaluates geometry for satellite positions
that transit near the polar axis.

\paragraph{Local East--North--Up frame.}
Many of the downstream models --- the tropospheric mapping function, the
elevation mask, the antenna phase-centre projection --- are most natural
in the topocentric ENU frame anchored at the receiver's geodetic
$(\geolat,\geolon)$.  The columns of the ENU$\to$ECEF rotation are the
three local unit vectors expressed in ECEF
(\texttt{Coordinates::enuToEcefMatrix}):
\begin{empheq}[box=\fbox]{equation}
\enuR(\geolat,\geolon) =
\big[\,\enuE \;\; \enuN \;\; \enuU\,\big]
=
\begin{bmatrix}
-\sin\geolon & -\sin\geolat\cos\geolon & \cos\geolat\cos\geolon \\
\;\;\cos\geolon & -\sin\geolat\sin\geolon & \cos\geolat\sin\geolon \\
0            & \cos\geolat              & \sin\geolat
\end{bmatrix}.
\label{eq:frame-enu}
\end{empheq}
The Up column is the outward geodetic normal --- identical to the
direction along which $\geoht$ is measured in \eqref{eq:frame-llh2ecef}
--- East is its horizontal cross-product with the pole, and North
completes the right-handed triad.  Because the columns are orthonormal,
$\enuR$ is a proper rotation and its inverse is its transpose: an ECEF
displacement is resolved into local components by
$\bm{v}_{\mathrm{ENU}} = \enuR^{\Tr}\,\bm{v}_{\mathrm{ECEF}}$
(\texttt{Coordinates::ecefToEnu}, with \texttt{Coordinates::enuToEcef}
the forward direction), avoiding any explicit matrix inversion.

\paragraph{Elevation and azimuth.}
With the receiver-to-satellite line-of-sight resolved into the local
frame, $\bm{e}_{\mathrm{ENU}} = (e_{\mathrm E}, e_{\mathrm N},
e_{\mathrm U})^{\Tr} = \enuR^{\Tr}\los$, the elevation and azimuth that
gate the measurement (the elevation mask of
\dcref{sec:obs-geometry}{the Observation Equations document} and the
mapping-function argument of \dcref{sec:atmos}{the Atmosphere document})
follow directly:
\begin{equation}
\elev = \arcsin\!\big(e_{\mathrm U}\big),
\qquad
\azim = \operatorname{atan2}\!\big(e_{\mathrm E},\, e_{\mathrm N}\big),
\label{eq:frame-elaz}
\end{equation}
azimuth measured clockwise from North.  Equations
\eqref{eq:frame-enu}--\eqref{eq:frame-elaz} are the entire geometric
bridge between the ECEF state vector and every elevation-dependent
weighting and antenna model in the reference.

% ============================================================
\subsection{Station reference points: marker, ARP, and phase centre}
\label{sec:frame-arp}
% ============================================================

Nothing above fixes which \emph{point} of a station a coordinate refers
to, and three distinct points are in play.  The \emph{marker}
(monument) is the surveyed benchmark, and is what an IGS SINEX/SSC
solution, an NGS Station Information File, and the RINEX
\texttt{APPROX POSITION XYZ} header record all refer to.  The
\emph{antenna reference point} (ARP) sits where the antenna meets its
mount and is the origin against which the ANTEX phase-centre offsets of
\dcref{sec:obsphys}{the Observation Physics document} are tabulated.
The \emph{antenna phase centre} is where the measurement is physically
realised.  The estimator's position state is the ARP --- for the simple
reason that the receiver PCO/PCV it applies are ARP-referenced --- so
the ARP$\to$phase-centre leg is closed by the antenna model and only the
marker$\to$ARP leg belongs to this document.  Confusing the two ends of
that leg puts a constant bias, equal to the eccentricity, straight into
the vertical.

The leg itself is the RINEX \texttt{ANTENNA: DELTA H/E/N} eccentricity
$\eccentricityHEN = (\delta H,\, \delta E,\, \delta N)^{\Tr}$, defined
as the offset from the marker to the ARP in the topocentric frame
\emph{at the marker}.  Note the ordering trap: RINEX writes
$(H,E,N) = (\text{Up},\text{East},\text{North})$, which is not the
$(E,N,U)$ order the rotation of \eqref{eq:frame-enu} consumes, so the
triple must be permuted before it is rotated
(\texttt{Coordinates::markerToArp}, the single owner of this
convention):
\begin{empheq}[box=\fbox]{equation}
\rxpos_{\mathrm{ARP}} \;=\; \rxpos_{\mathrm{mkr}}
  \;+\; \enuR\big(\geolat_{\mathrm{mkr}},\, \geolon_{\mathrm{mkr}}\big)
    \begin{bmatrix}\delta E\\[2pt] \delta N\\[2pt] \delta H\end{bmatrix},
\label{eq:frame-marker2arp}
\end{empheq}
with $(\geolat_{\mathrm{mkr}},\geolon_{\mathrm{mkr}})$ recovered from
$\rxpos_{\mathrm{mkr}}$ by the inverse map of
\eqref{eq:frame-bowring-iter}.  Write $\eccentricityHEN'$ for the
permuted $(\delta E, \delta N, \delta H)^{\Tr}$ that the rotation acts
on.

\paragraph{The inverse is not a sign flip.}
Reducing an ARP-referenced estimate back to the monument --- to compare
it against, or publish it in, a surveyed coordinate --- cannot simply
negate \eqref{eq:frame-marker2arp}, because the local frame belongs at
the marker and the marker is precisely the unknown.  Evaluating $\enuR$
at the ARP instead leaves a closure error of order
$\lVert\eccentricityHEN\rVert^{2}/R_{\oplus}$.  The engine therefore
takes one fixed-point step (\texttt{Coordinates::arpToMarker}):
\begin{equation}
\rxpos^{(0)} = \rxpos_{\mathrm{ARP}}
  - \enuR\big(\mathrm{llh}(\rxpos_{\mathrm{ARP}})\big)\,\eccentricityHEN',
\qquad
\rxpos_{\mathrm{mkr}} = \rxpos_{\mathrm{ARP}}
  - \enuR\big(\mathrm{llh}(\rxpos^{(0)})\big)\,\eccentricityHEN',
\label{eq:frame-arp2marker}
\end{equation}
which closes the round trip with \eqref{eq:frame-marker2arp} to well
below the millimetre level this reference grades at.

\paragraph{Where the reduction is applied --- and where it is not.}
Every position derived from a RINEX \emph{header} is reduced: the
post-processing filter seed (\texttt{Application::initializePosition}
and the \texttt{PppModel} state initialiser), the real-time seed
(\texttt{RealTimeEngine::initialize}), the baseline rover a~priori, and
the closed-loop consumer truth all take the marker through
\eqref{eq:frame-marker2arp} --- some calling
\texttt{Coordinates::markerToArp} directly, others through the accessor
\texttt{RinexParser::getArpPositionXyz} that pairs the two header
fields.  A coordinate supplied
\emph{directly in the configuration} is not reduced: it is taken at face
value, on the stated policy that the user owns its reference point.  So
the baseline \emph{reference} coordinate is used as given while the
rover a~priori read from its header is reduced --- an asymmetry inside
one mode --- and, more consequentially, the network producer's station
coordinates (\texttt{StationConfig::position\_ecef}, optionally
SIF-resolved under \texttt{network.use\_sif\_coordinates}) are used as
given while being held \emph{fixed}, their three position columns
dropped from the normal equations.  A monument-referenced coordinate
there is therefore not absorbed by a position state; it survives into
the residuals and hence into the estimated clocks, biases, and
ionosphere.  A zero or absent eccentricity makes
\eqref{eq:frame-marker2arp} an exact, bitwise no-op, which is what keeps
header-less RTCM/SBF streams and \texttt{0/0/0} stations untouched; the
ordering, sign, no-op, and round-trip properties are pinned by the
\texttt{AntennaEccentricityTest} suite.

% ============================================================
\subsection{GNSS time scales and the GPS reference epoch}
\label{sec:frame-timescale}
% ============================================================

Time inside the engine is carried as a \texttt{GnssTime}
$\{\text{week},\text{tow}\}$ pair --- an integer GPS week and a
floating-point time-of-week in seconds --- and is expanded to a single
continuous count $T = \text{week}\times 604800 + \text{tow}$ wherever a
monotone scalar is needed (e.g.\ ERP interpolation,
\Cref{sec:frame-erp}).  A week boundary is $604\,800\;\mathrm{s}$ and a
day $86\,400\;\mathrm{s}$ (\texttt{src/core/math/TimeUtils.cpp}, header
constants \texttt{SECONDS\_IN\_WEEK}, \texttt{SECONDS\_IN\_DAY}).  The
common origin is the GPS epoch, 1980-01-06 00:00:00 UTC, whose Julian
Date is fixed in the code:
\begin{equation}
\mathrm{JD}_{\mathrm{GPS}} = 2444244.5
\qquad(\texttt{src/core/math/TimeUtils.cpp}, \texttt{GPS\_EPOCH\_JD}).
\label{eq:frame-gpsepoch}
\end{equation}
GPS time is a continuous atomic scale --- it never inserts a leap second
--- so it sits a fixed integer offset below International Atomic Time and
a fixed fractional offset below Terrestrial Time
(\texttt{src/core/math/TimeUtils.cpp}, header constants
\texttt{TAI\_GPS\_OFFSET}, \texttt{TT\_GPS\_OFFSET}):
\begin{empheq}[box=\fbox]{equation}
\mathrm{TAI} - \mathrm{GPS} = 19\;\mathrm{s},
\qquad
\mathrm{TT} - \mathrm{GPS} = 51.184\;\mathrm{s} = (32.184 + 19)\;\mathrm{s}.
\label{eq:frame-tai-tt}
\end{empheq}
The $32.184\;\mathrm s$ in the TT relation is the constant TT$-$TAI
offset; together these let any GPS time be promoted to the TT scale on
which precession and nutation are evaluated.  Terrestrial Time relative
to the J2000.0 epoch is measured in Julian centuries,
\begin{equation}
t = \frac{\mathrm{JD}_{\mathrm{TT}} - 2451545.0}{36525.0},
\label{eq:frame-jcent}
\end{equation}
the dimensionless argument of the astronomical series of
\Cref{sec:frame-eci,sec:frame-sunmoon}
(\texttt{TimeUtils::jdToJ2000Centuries}, constants \texttt{J2000\_JD},
\texttt{DAYS\_PER\_JULIAN\_CENTURY}).

\paragraph{Gregorian$\leftrightarrow$Julian Date.}
The calendar conversions follow the Meeus algorithm \cite{meeus1991}.
The forward map, with the month shifted so January and February count as
months 13 and 14 of the prior year and a Gregorian-reform correction
$b = 2 - A + \lfloor A/4\rfloor$ for $A = \lfloor y/100\rfloor$, is
(\texttt{TimeUtils::gregorianToJulianDay}):
\begin{equation}
\mathrm{JD} = \big\lfloor 365.25\,(y+4716)\big\rfloor
            + \big\lfloor 30.6001\,(m+1)\big\rfloor
            + d + b - 1524.5 + \frac{\text{(time of day)}}{86400},
\label{eq:frame-greg2jd}
\end{equation}
and the inverse (\texttt{TimeUtils::julianDayToGregorian}) reverses
it through the integer auxiliaries $z = \lfloor\mathrm{JD}+0.5\rfloor$
and the Gregorian-reform branch at $z \ge 2299161$.  The crucial
distinction the code maintains throughout is between a calendar date that
is already on the GPS scale (no leap correction:
\texttt{gregorianGpsToGpsTime}) and one that is on UTC and must have the
epoch's leap seconds added back to land on GPS time
(\texttt{TimeUtils::gregorianToGpsTime}); conflating the two is the
classic source of whole-second timing errors.

% ============================================================
\subsection{Leap seconds and per-system dispatch}
\label{sec:frame-leap}
% ============================================================

UTC is tied to GPS by the historical table of leap-second insertions.
The engine stores the table as $\big(\mathrm{JD}_k,\,\leapgu(k)\big)$
pairs --- the GPS$-$UTC offset that takes effect at each insertion
boundary, the first row being the GPS epoch itself with offset zero
(\texttt{src/core/math/TimeUtils.cpp}, \texttt{LEAP\_TABLE}).  A query resolves the
offset by a reverse scan that returns the offset of the latest boundary
not after the epoch:
\begin{equation}
\leapgu(\mathrm{JD}) = \leapgu(k^{\star}),
\qquad
k^{\star} = \max\{\,k : \mathrm{JD}_k \le \mathrm{JD}\,\},
\label{eq:frame-leaplookup}
\end{equation}
falling back to zero before the epoch
(\texttt{TimeUtils::leapSecondsGpsUtc}).  TAI$-$UTC is then simply
$\leapgu + 19\;\mathrm{s}$ from \eqref{eq:frame-tai-tt}
(\texttt{TimeUtils::leapSecondsTaiUtc}).  The epoch-dependence is
essential: the same UTC calendar instant maps to different GPS times in
2005 and 2010 because a leap second was inserted between them, and only a
boundary-aware lookup recovers the correct offset.

\paragraph{The wall-clock leap constant and its runtime override.}
The boundary-aware table of \eqref{eq:frame-leaplookup} is the authority
for every \emph{epoch-referenced} conversion --- a dated calendar
timestamp, GLONASS day arithmetic, UTC formatting --- because each of
those carries the epoch at which the correct offset is known.  The one
conversion that does not is the promotion of the \emph{system wall clock}
to GPS time (\texttt{TimeUtils::nowContinuousGpsSeconds},
\texttt{nowGpsTime}), used only on the real-time decode and
latency-stamping paths: POSIX \texttt{system\_clock} is UTC and carries no
epoch of its own, so the offset cannot be looked up and must instead be
supplied as a fixed count.  That count is the compiled constant
\texttt{TimeUtils::LEAP\_SECONDS} $= 18\;\mathrm{s}$ (valid from
2017-01-01), read through an atomic accessor
(\texttt{currentLeapSeconds}) so that a future insertion can be applied to
a live binary without a rebuild: \texttt{setLeapSecondsOverride}, driven
by the config key \texttt{realtime.leap\_seconds}
(\texttt{src/core/config/GnssConfigLoad.cpp}), replaces it at real-time
startup.  The override deliberately touches \emph{only} the wall-clock
path; the epoch-dependent \eqref{eq:frame-leaplookup} lookups are left
alone, since they are correct at every historical epoch outside the
insertion instants themselves\footnote{The comparison in
\texttt{TimeUtils::leapSecondsGpsUtc} is made on the GPS scale against
boundary Julian Dates tabulated at the \emph{UTC} insertion instant, so
the offset advances $\leapgu$ seconds early --- a known, deferred
discrepancy confined to a $\leapgu$-second window at each boundary
(18\,s at the most recent insertion, 2017-01-01).} and an override
there would corrupt rather than update them.

\paragraph{Per-constellation Gregorian dispatch.}
The four supported time systems differ only in how a RINEX/broadcast
calendar timestamp is referred to the GPS scale, dispatched by a system
tag (\texttt{TimeUtils::GnssTimeFromDate}):
\begin{itemize}
  \item \textbf{GPS / Galileo} --- already on the GPS scale; the
    timestamp is converted with no leap correction (GST is steered to
    GPS time and shares the same week/second framing to within the few
    nanoseconds the ISB state of \dcref{sec:state}{the State Vector
    document} absorbs).
  \item \textbf{BeiDou} --- the timestamp is on BDT, which runs
    $14\;\mathrm{s}$ behind GPS (\Cref{sec:frame-bdt}); the code converts
    on the GPS-scale assumption and then \emph{adds} $14\;\mathrm{s}$,
    carrying any week roll-over
    (\texttt{TimeUtils::GnssTimeFromDate}, BeiDou branch).
  \item \textbf{GLONASS} --- the RINEX timestamp is UTC; the full leap
    correction of \eqref{eq:frame-leaplookup} is applied to bring it onto
    GPS time (\texttt{TimeUtils::GnssTimeFromDate}, GLONASS branch).
\end{itemize}
GPS and Galileo therefore carry no leap correction, BeiDou a fixed
$14\;\mathrm{s}$, and GLONASS the full epoch-dependent leap count --- a
dispatch that exactly mirrors the constellations' defining time
references.  There is a fourth branch: an unrecognised time-system tag
falls through to the GPS-scale conversion with only a trace-level note,
and a calendar date that fails validation returns the sentinel
$\{0,0\}$ rather than a fabricated epoch
(\texttt{TimeUtils::GnssTimeFromDate}).  The fallback is silent at
default verbosity, so an unexpected tag mis-times by the full leap count
rather than failing loudly.

% ============================================================
\subsection{GLONASS time-of-day to full GPS week}
\label{sec:frame-glonass}
% ============================================================

GLONASS broadcast and SSR records frequently carry only a time-of-day
--- seconds since the start of the GLONASS day --- with no week and no
absolute date.  GLONASS time runs on UTC(SU) advanced by three hours
(Moscow), so its day boundary is three hours ahead of the UTC midnight.
Reconstructing a full GPS week from a bare time-of-day therefore needs an
\emph{anchor}: an approximate GPS-seconds estimate of the true epoch
(typically the current obs epoch), against which the day and the
$\pm12\,\mathrm{h}$ wrap are resolved
(\texttt{TimeUtils::glonassTodToGpsTime}, after RTKLIB's
\texttt{adjday\_glot}).  The procedure, with day length
$D = 86\,400\;\mathrm{s}$, half-day $D/2$, and Moscow offset
$M = 3\times3600\;\mathrm{s}$:
\begin{enumerate}
  \item Convert the anchor GPS seconds to UTC by subtracting the leap
    count at that epoch, $\,t_{\mathrm{UTC}} = t_{\mathrm{anchor}} -
    \leapgu$, then shift into the Moscow-day frame
    $t_{\mathrm M} = t_{\mathrm{UTC}} + M$ and locate that day's midnight
    $t_0 = \lfloor t_{\mathrm M}/D\rfloor\,D$, giving the anchor's
    time-of-day $\tau_{\mathrm a} = t_{\mathrm M} - t_0$.
  \item Place the reported time-of-day $\tau$ in the anchor's day,
    resolving the midnight wrap toward the nearest match:
\end{enumerate}
\begin{empheq}[box=\fbox]{equation}
\tau \leftarrow
\begin{cases}
\tau + D, & \tau < \tau_{\mathrm a} - D/2,\\[2pt]
\tau - D, & \tau > \tau_{\mathrm a} + D/2,\\[2pt]
\tau,     & \text{otherwise,}
\end{cases}
\qquad
t_{\mathrm{GPS}} = \underbrace{(t_0 + \tau)}_{\text{Moscow}} - M + \leapgu .
\label{eq:frame-glonass-tod}
\end{empheq}
The final line undoes the Moscow offset to recover UTC and then
\emph{re-adds the same leap count} $\leapgu$ to climb back onto the GPS
scale.  The leap term enters once with each sign and so cancels exactly
between the two legs --- it is present only because the intermediate
arithmetic must pass through UTC to frame the Moscow day correctly.  If
no usable anchor is supplied (a non-positive value) the routine returns
the sentinel $\{0,0\}$ rather than fabricate a week, so the caller can
detect the failure and fall back; the continuous-seconds-to-week split
that closes the routine carries the standard negative-TOW correction
(\texttt{gpsSecondsToGnssTime}).

% ============================================================
\subsection{BeiDou time: the 14\,s offset and independent week base}
\label{sec:frame-bdt}
% ============================================================

BeiDou time (BDT) is offset from GPS by a sub-second-level integer and by
an independent week numbering, and the two pieces are owned by different
modules.  The \emph{sub-second} relation,
\begin{equation}
\mathrm{BDT} = \mathrm{GPS} - 14\;\mathrm{s},
\label{eq:frame-bdt-sec}
\end{equation}
reflects that BDT was synchronised to UTC at its 2006 zero epoch, by
which time GPS already led UTC by $14\;\mathrm{s}$; it is applied in
\texttt{TimeUtils} when a BeiDou calendar timestamp is converted to GPS
time (\texttt{TimeUtils::GnssTimeFromDate}, BeiDou branch; \Cref{sec:frame-leap}).
The \emph{week} relation is separate.  BDT counts weeks from its own zero
epoch (2006-01-01), which is GPS week 1356, and this base is applied at
query time by the ephemeris store, not by \texttt{TimeUtils}
(the file-local constant \texttt{kBdtToGpsOffsetSec}, consumed by
\texttt{LiveEphemerisStore::computeSatellitePV}).  When the
store lifts a BeiDou ephemeris reference time onto the common GPS scale
for an honest cross-system age check it adds the combined offset
\begin{empheq}[box=\fbox]{equation}
\Delta_{\mathrm{BDT}\to\mathrm{GPS}}
= 1356 \times 604800\;\mathrm{s} + 14\;\mathrm{s},
\label{eq:frame-bdt-week}
\end{empheq}
the week base plus the $14\;\mathrm{s}$ of \eqref{eq:frame-bdt-sec}
(\texttt{kBdtToGpsOffsetSec}, applied by the file-local
\texttt{toeGpsContinuous}).  Omitting the
week base would make every BeiDou satellite read roughly 1356 weeks
($\sim$26 years) stale and be silently rejected by the staleness guard
that follows; the offset does not cancel automatically because the query
time is GPS-scale while the stored $t_{\mathrm{oe}}$ is on the native BDT
week (the staleness guard in
\texttt{LiveEphemerisStore::computeSatellitePV}).

% ============================================================
\subsection{Legacy ECI$\to$ECEF: the equinox-based chain}
\label{sec:frame-eci}
% ============================================================

The Sun, Moon, and (for some product streams) the satellite positions
are most naturally given in a quasi-inertial celestial frame; the tides,
attitude, and phase wind-up that consume them live in the terrestrial
ECEF frame.  The bridge is the Earth-orientation matrix $\Umat$ such that
$\bm{r}_{\mathrm{ECEF}} = \Umat\,\bm{r}_{\mathrm{ECI}}$.  The default
(always-available) path is the classical equinox-based chain following
the IAU-1976 precession and IAU-1980 nutation conventions
\cite{iers2010}, composed as (\texttt{Planetary::eci2ecef}):
\begin{empheq}[box=\fbox]{equation}
\Umat = \polW\;\rotR{3}(\gast)\;\nutN\;\precP .
\label{eq:frame-legacy-chain}
\end{empheq}
Here $\rotR{1},\rotR{2},\rotR{3}$ are the elementary right-handed axis
rotations of \texttt{Planetary::Rxyz} and
each factor reverses one piece of the celestial-to-terrestrial
orientation: precession and nutation orient the true equator and equinox
of date, the sidereal rotation $\rotR{3}(\gast)$ spins the frame about
the pole by the apparent sidereal time, and polar motion $\polW$ ties the
celestial pole to the conventional terrestrial pole.

\paragraph{Precession (IAU 1976).}
With $t$ the TT centuries of \eqref{eq:frame-jcent} and arcsecond
coefficients converted to radians, the three equatorial precession angles
$\zeta_A, \theta_A, z_A$ and the mean obliquity $\obleps_0$ are
(the precession block of \texttt{Planetary::eci2ecef}):
\begin{align}
\zeta_A &= 2306.2181\,t + 0.30188\,t^2 + 0.017998\,t^3, \nonumber\\
\theta_A &= 2004.3109\,t - 0.42665\,t^2 - 0.041833\,t^3, \nonumber\\
z_A &= 2306.2181\,t + 1.09468\,t^2 + 0.018203\,t^3,
\label{eq:frame-prec-angles}\\
\obleps_0 &= 84381.448 - 46.8150\,t - 0.00059\,t^2 + 0.001813\,t^3
\quad(\text{arcsec}), \nonumber
\end{align}
assembled as
$\precP = \rotR{3}(-z_A)\,\rotR{2}(\theta_A)\,\rotR{3}(-\zeta_A)$.

\paragraph{Nutation (IAU 1980, 106 terms).}
The nutation in longitude $\dpsi$ and in obliquity $\deps$ are summed
over the 106-term IAU-1980 series \cite{iers2010} in the five Delaunay
fundamental arguments $f = (l, l', F, D, \Omega)$
(\texttt{Planetary::nutIau1980}; arguments from
\texttt{Planetary::astArgs}):
\begin{equation}
\dpsi = 10^{-4}\!\!\sum_{i=1}^{106}\!\big(A^{\psi}_i + \dot A^{\psi}_i\,t\big)\sin\!\big(\bm{n}_i\!\cdot\! f\big),
\quad
\deps = 10^{-4}\!\!\sum_{i=1}^{106}\!\big(A^{\varepsilon}_i + \dot A^{\varepsilon}_i\,t\big)\cos\!\big(\bm{n}_i\!\cdot\! f\big),
\label{eq:frame-nutation}
\end{equation}
in arcseconds (the $10^{-4}$ scales the tabulated $0.1\,\mathrm{mas}$
units), then converted to radians.  The nutation matrix rotates by the
true obliquity $\obleps_0 + \deps$:
$\nutN = \rotR{1}(-\obleps_0-\deps)\,\rotR{3}(-\dpsi)\,\rotR{1}(\obleps_0)$
(the nutation block of \texttt{Planetary::eci2ecef}).

\paragraph{Apparent sidereal time and the 0h-UT1 subtlety.}
Greenwich Apparent Sidereal Time adds the equation of the equinoxes ---
the nutation in longitude projected on the equator, plus the two leading
complementary terms in the lunar node $\Omega$ --- to the mean sidereal
time (the GAST block of \texttt{Planetary::eci2ecef}):
\begin{equation}
\gast = \gmst + \dpsi\cos\obleps_0
        + \big(0.00264\sin\Omega + 0.000063\sin 2\Omega\big)\,\text{arcsec}.
\label{eq:frame-gast}
\end{equation}
GMST itself is built from the IAU-1982 model evaluated at the
\emph{start} of the current UT1 day plus the Earth's rotation through the
elapsed day fraction (\texttt{Planetary::utc2gmst}).  The reference
instant must be 0h UT1 of the calendar day, and because a Julian Date is
referenced to \emph{noon}, that midnight is recovered as
\begin{equation}
\mathrm{JD}_{0\mathrm h} = \big\lfloor \mathrm{JD}_{\mathrm{UT1}} - 0.5\big\rfloor + 0.5,
\label{eq:frame-0h-ut1}
\end{equation}
not $\lfloor\mathrm{JD}\rfloor$ (same routine).  Using
the bare floor lands on the \emph{previous} noon for any post-midnight
epoch --- a half-day error that propagates to roughly $0.5^{\circ}$ in
GMST and, through \eqref{eq:frame-legacy-chain}, on the order of
$10^{6}\;\mathrm{km}$ in the rotated Sun position; the comment in
\texttt{Planetary::utc2gmst} records exactly this failure mode.

\paragraph{Polar motion.}
The terrestrial-pole tie is
$\polW = \rotR{2}(-x_p)\,\rotR{1}(-y_p)$ for the interpolated pole
coordinates $(x_p, y_p)$ in radians
(the polar-motion block of \texttt{Planetary::eci2ecef}), the small-angle
two-axis rotation that aligns the celestial intermediate pole of date
with the IERS conventional pole.

% ============================================================
\subsection{CIO-based ITRS$\leftrightarrow$GCRS (IAU 2000A/2006)}
\label{sec:frame-cio}
% ============================================================

An opt-in path (config key \texttt{estimation.enable\_cio\_frame}, getter
\texttt{GnssConfig::enableCioFrame}; \textbf{default off}, so the
equinox chain of \eqref{eq:frame-legacy-chain} is the shipped path)
replaces the equinox chain
with the modern Celestial Intermediate Origin (CIO) based transform of
the IAU 2000A/2006 conventions \cite{iers2010}, evaluated through SOFA
\cite{sofa} and mirroring the M-PAGES \texttt{NgsFormC2T} construction
(\texttt{Planetary::celestialToTerrestrial}).  The
celestial-to-terrestrial matrix is composed as
\begin{empheq}[box=\fbox]{equation}
\Umat = \mathtt{POM}(x_p, y_p, s')\;\rotR{3}(\era)\;\mathtt{C2I}(\cipX,\cipY,\cioloc),
\label{eq:frame-cio-chain}
\end{empheq}
i.e.\ the SOFA composition \texttt{iauC2tcio}.  The Celestial
Intermediate Pole coordinates $(\cipX,\cipY)$ and the CIO locator
$\cioloc$ come from the full IAU-2000A precession-nutation model via
\texttt{iauXys00a}, after which the \emph{observed} celestial-pole
offsets are injected additively,
\begin{equation}
\cipX \leftarrow \cipX + \mathrm{d}X, \qquad
\cipY \leftarrow \cipY + \mathrm{d}Y,
\label{eq:frame-cip-offset}
\end{equation}
and the celestial-to-intermediate matrix is formed by \texttt{iauC2ixys}.
The diurnal rotation is the Earth Rotation Angle $\era$ (a strictly
linear function of UT1, \texttt{iauEra00}) rather than the sidereal time
of \eqref{eq:frame-gast}, and the terrestrial-pole tie
$\mathtt{POM}$ uses \texttt{iauPom00} with the TIO locator $s'$ from
\texttt{iauSp00}.  Two-part Julian Dates are passed with the
$(2400000.5, \mathrm{MJD})$ split for numerical conditioning
(same routine).  The two paths are designed to agree:
with $\mathrm{d}X = \mathrm{d}Y = 0$ the CIO result reproduces the bare
IAU-2000A model and differs from the legacy equinox chain only by the
(sub-milliarcsecond) frame-bias and model-version difference, which the
engine confirms at startup by logging the angular separation of the two
Sun directions (the startup diagnostic in
\texttt{Planetary::computeSunMoonPosition}).

% ============================================================
\subsection{Earth-orientation parameters: parse and interpolation}
\label{sec:frame-erp}
% ============================================================

Both transforms above need the time-varying Earth-orientation parameters
$(x_p, y_p, \mathrm{UT1}-\mathrm{UTC}, \mathrm{d}X, \mathrm{d}Y)$.  These
arrive in a standard IGS \texttt{.ERP} file, parsed by
\texttt{ProductParser::parseErp} (\texttt{src/core/products/ProductParser.cpp}).
The parser reads the MJD time tag and the pole/UT1 columns and applies
the IGS scale factors --- the pole coordinates are tabulated in
$10^{-6}$ arcseconds and UT1$-$UTC in $10^{-7}$ seconds
(\texttt{ProductParser::parseErp}):
\begin{equation}
x_p = (\text{col})\times 10^{-6}\,\text{arcsec}, \quad
y_p = (\text{col})\times 10^{-6}\,\text{arcsec}, \quad
(\mathrm{UT1}-\mathrm{UTC}) = (\text{col})\times 10^{-7}\,\text{s}.
\label{eq:frame-erp-scale}
\end{equation}
The MJD tag is converted to a GPS week/TOW relative to the GPS-epoch MJD
(same routine).  The CIP offsets $\mathrm{d}X,
\mathrm{d}Y$ are \emph{not} carried by the standard \texttt{.ERP} format
and default to zero; because the IAU-2000A model of
\eqref{eq:frame-cio-chain} already includes the modelled
precession-nutation, $\mathrm{d}X = \mathrm{d}Y = 0$ is the correct
nominal value, not a missing input (same routine).  The parser
\emph{only} parses;
all temporal interpolation lives in \texttt{Planetary::interpolateErp},
which binary-searches
the chronologically sorted table on continuous GPS seconds and linearly
interpolates every component between the bracketing entries:
\begin{equation}
p(t) = p_{0} + \frac{t - t_{0}}{t_{1}-t_{0}}\,(p_{1}-p_{0}),
\qquad p \in \{x_p, y_p, \mathrm{UT1}{-}\mathrm{UTC}, \mathrm{d}X, \mathrm{d}Y\},
\label{eq:frame-erp-interp}
\end{equation}
clamping to the nearest endpoint outside the table span (same routine).
The interpolated $(x_p, y_p)$ feed
the polar-motion factor of \eqref{eq:frame-legacy-chain} /
\eqref{eq:frame-cio-chain}, while UT1$-$UTC corrects UTC to UT1 before
the sidereal/ERA evaluation.  The product-evaluation document
(\dcref{sec:products}{Product Evaluation}) cross-references this section
rather than re-deriving the polar-motion application, since the same
$\polW$ rotates both the luni-solar vectors and any celestial-frame
satellite product.

% ============================================================
\subsection{The Sun and Moon series these transforms position}
\label{sec:frame-sunmoon}
% ============================================================

The end purpose of the orientation chain is to place the Sun and Moon in
ECEF for the geophysical and observation-model corrections.  The
celestial-frame positions are computed first, from an analytic series in
the J2000 centuries $t$ and the Delaunay arguments $f$
(\texttt{Planetary::sunMoonPosEci}), an enhanced
RTKLIB/Meeus model \cite{meeus1991, montenbruck2000}.  The Sun's
ecliptic longitude is the mean longitude plus the equation of centre
(three harmonics of the solar mean anomaly $M_\odot$) and small Venus,
Jupiter, and lunar perturbation terms (same routine); the obliquity
rotation
\begin{equation}
\obleps = 23.439291 - 0.0130042\,t \quad(\text{deg})
\label{eq:frame-obliquity}
\end{equation}
carries the resulting ecliptic
$(\lambda_\odot, \beta_\odot)$ into equatorial ECI components:
\begin{align}
r_x &= r\,\cos\beta_\odot\cos\lambda_\odot, \nonumber\\
r_y &= r\,(\cos\obleps\,\cos\beta_\odot\sin\lambda_\odot - \sin\obleps\,\sin\beta_\odot), \label{eq:frame-eci-sun}\\
r_z &= r\,(\sin\obleps\,\cos\beta_\odot\sin\lambda_\odot + \cos\obleps\,\sin\beta_\odot), \nonumber
\end{align}
with the solar distance $r$ from the three-harmonic radius series.  The
Moon's longitude, latitude, and horizontal-parallax-derived distance
follow the analogous Brown-series truncation in the same five
fundamental arguments (same routine).  The resulting ECI vectors are
then
rotated to ECEF by exactly the $\Umat$ of \eqref{eq:frame-legacy-chain}
or \eqref{eq:frame-cio-chain} according to the
\texttt{estimation.enable\_cio\_frame} toggle (default off,
\Cref{sec:frame-cio}), in
\texttt{Planetary::computeSunMoonPosition}; the ECEF Sun and Moon are precisely
the inputs consumed by the solid-Earth and ocean tides of
\dcref{sec:sitedisp}{the Site Displacement document} and by the
satellite-attitude and phase-wind-up models of
\dcref{sec:obsphys}{the Observation Physics document}.  The two
luni-solar models are only $\sim$100\,km accurate in position --- ample
for the tidal and attitude geometry, which depend on the Sun/Moon
\emph{direction} to well under the milliarcsecond level that matters ---
and a non-finite guard substitutes a nominal 1\,AU / mean-lunar-distance
fallback rather than propagate a NaN into the correction stack (the
output-validation tail of \texttt{Planetary::computeSunMoonPosition}).

\ifSubfilesClassLoaded{\printbibliography}{}
\end{document}
